% GENERATED FILE. Edit canonical lessons, then run npm run generate:books.
% canonical-source-hash: fnv1a64:6fdbd76aae80c5ea
% canonical-lessons: FA-C326-masdum, FA-C326-jerahat, FA-C326-khunrizi, FA-C326-bihush-shodan, FA-C326-be-hush-amadan

\chapter{The injured}
\label{ch:fa-the-injured}
\hlchaptermodality{326}

\begin{chapteropening}
\textbf{By the end of this chapter:} \emph{I can talk about the injured: a wound, bleeding, fainting and coming round.}

Complete the last lesson of chapter 326: I can talk about the injured: a wound, bleeding, fainting and coming round.
\end{chapteropening}

\section[\texorpdfstring{\fa{مصدوم} (masdum)}{مصدوم (masdum)}]{\fa{مصدوم} (masdum) — injured; an injured person}

\label{lesson:FA-C326-masdum}



\begin{quote}
Before the new one: say the Persian for to be harmed, then the Persian for a fracture.
\end{quote}

\subsection*{You'll want to know: \fa{مصدوم}}
\textbf{\fa{مصدوم}} (\emph{masdum}) — \textquotedblleft{}injured; an injured person\textquotedblright{}.

\begin{quote}
\fa{مصدوم} \fa{را} \fa{با} \fa{آمبولانس} \fa{به} \fa{بیمارستان} \fa{بردند}.
\end{quote}

\begin{quote}
\emph{Masdum râ bâ âmbulâns be bimârestân bordand.}
\end{quote}

\begin{quote}
They took the injured man to hospital by ambulance.
\end{quote}

In news reports; plural \emph{masdumân} or \emph{masdumin}.

\begin{cousinweb}
\emph{masdum} is usually traced to Arabic \emph{maṣdūm}, \textquotedblleft{}struck\textquotedblright{}, a participle of the same root as \emph{ṣadma} \textquotedblleft{}a blow\textquotedblright{}; EtymDB lists no entry.
\end{cousinweb}

\begin{grammarlens}[title={What you've built}]
One more word for telling, explaining and describing at B1.
\end{grammarlens}

\subsection*{Your turn}
\begin{itemize}
\raggedright
  \item \emph{Say it:} \emph{masdum}
  \item \emph{Say it:} \emph{masdum}, once more
  \item \emph{Say it:} a sentence of your own with \emph{masdum}
  \item \emph{Recall:} say the Persian for to be harmed, then the Persian for a fracture, then say \emph{masdum} again
\end{itemize}

\begin{tcolorbox}[breakable,colback=teal!4,colframe=teal!35!black,title={Before you move on}]
What does \fa{مصدوم} mean? (\textquotedblleft{}Injured; an injured person\textquotedblright{}.) Say it once more.
\end{tcolorbox}
\section[\texorpdfstring{\fa{جراحت} (jerâhat)}{جراحت (jerâhat)}]{\fa{جراحت} (jerâhat) — a wound, an injury}

\label{lesson:FA-C326-jerahat}



\begin{quote}
Before the new one: say the Persian for a fracture, then the Persian for injured.
\end{quote}

\subsection*{You'll want to know: \fa{جراحت}}
\textbf{\fa{جراحت}} (\emph{jerâhat}) — \textquotedblleft{}a wound, an injury\textquotedblright{}.

\begin{quote}
\fa{جراحت} \fa{او} \fa{عمیق} \fa{نبود}.
\end{quote}

\begin{quote}
\emph{Jerâhat-e u 'amiq nabud.}
\end{quote}

\begin{quote}
His wound was not deep.
\end{quote}

Formal; in speech \emph{zakhm}. The same root gives \emph{jarrâh} \textquotedblleft{}a surgeon\textquotedblright{}.

\begin{cousinweb}
\emph{jerâhat} is usually traced to Arabic \emph{jirāḥa}, \textquotedblleft{}wounds, wounding\textquotedblright{}; EtymDB lists no entry.
\end{cousinweb}

\begin{grammarlens}[title={What you've built}]
One more word for telling, explaining and describing at B1.
\end{grammarlens}

\subsection*{Your turn}
\begin{itemize}
\raggedright
  \item \emph{Say it:} \emph{jerâhat}
  \item \emph{Say it:} \emph{jerâhat}, once more
  \item \emph{Say it:} a sentence of your own with \emph{jerâhat}
  \item \emph{Recall:} say the Persian for a fracture, then the Persian for injured, then say \emph{jerâhat} again
\end{itemize}

\begin{tcolorbox}[breakable,colback=teal!4,colframe=teal!35!black,title={Before you move on}]
What does \fa{جراحت} mean? (\textquotedblleft{}A wound, an injury\textquotedblright{}.) Say it once more.
\end{tcolorbox}
\section[\texorpdfstring{\fa{خونریزی} (khunrizi)}{خونریزی (khunrizi)}]{\fa{خونریزی} (khunrizi) — bleeding; bloodshed}

\label{lesson:FA-C326-khunrizi}



\begin{quote}
Before the new one: say the Persian for injured, then the Persian for a wound.
\end{quote}

\subsection*{You'll want to know: \fa{خونریزی}}
\textbf{\fa{خونریزی}} (\emph{khunrizi}) — \textquotedblleft{}bleeding; bloodshed\textquotedblright{}.

\begin{quote}
\fa{پرستار} \fa{خونریزی} \fa{را} \fa{متوقف} \fa{کرد}.
\end{quote}

\begin{quote}
\emph{Parastâr khunrizi râ motevaqqef kard.}
\end{quote}

\begin{quote}
The nurse stopped the bleeding.
\end{quote}

\emph{khun} \textquotedblleft{}blood\textquotedblright{} + \emph{rizi} \textquotedblleft{}pouring\textquotedblright{}, from \emph{riz-}, the present stem of \emph{rikhtan}.

\begin{cousinweb}
\emph{khunrizi} joins \emph{khun} \textquotedblleft{}blood\textquotedblright{} and \emph{rizi} \textquotedblleft{}pouring\textquotedblright{}, from \emph{rikhtan} \textquotedblleft{}to pour\textquotedblright{}, which EtymDB traces to Middle Persian. All of it is native Persian.
\end{cousinweb}

\begin{grammarlens}[title={What you've built}]
One more word for telling, explaining and describing at B1.
\end{grammarlens}

\subsection*{Your turn}
\begin{itemize}
\raggedright
  \item \emph{Say it:} \emph{khunrizi}
  \item \emph{Say it:} \emph{khunrizi}, once more
  \item \emph{Say it:} a sentence of your own with \emph{khunrizi}
  \item \emph{Recall:} say the Persian for injured, then the Persian for a wound, then say \emph{khunrizi} again
\end{itemize}

\begin{tcolorbox}[breakable,colback=teal!4,colframe=teal!35!black,title={Before you move on}]
What does \fa{خونریزی} mean? (\textquotedblleft{}Bleeding; bloodshed\textquotedblright{}.) Say it once more.
\end{tcolorbox}
\section[\texorpdfstring{\fa{بیهوش} \fa{شدن} (bihush shodan)}{بیهوش شدن (bihush shodan)}]{\fa{بیهوش} \fa{شدن} (bihush shodan) — to faint, to pass out}

\label{lesson:FA-C326-bihush-shodan}



\begin{quote}
Before the new one: say the Persian for a wound, then the Persian for bleeding.
\end{quote}

\subsection*{You'll want to know: \fa{بیهوش} \fa{شدن}}
\textbf{\fa{بیهوش} \fa{شدن}} (\emph{bihush shodan}) — \textquotedblleft{}to faint, to pass out\textquotedblright{}.

\begin{quote}
\fa{زن} \fa{از} \fa{گرما} \fa{بیهوش} \fa{شد}.
\end{quote}

\begin{quote}
\emph{Zan az garmâ bihush shod.}
\end{quote}

\begin{quote}
The woman fainted from the heat.
\end{quote}

\emph{bi} \textquotedblleft{}without\textquotedblright{} + \emph{hush} \textquotedblleft{}consciousness\textquotedblright{}. Present \emph{bihush mishavad}.

\begin{cousinweb}
\emph{hush} \textquotedblleft{}consciousness, mind\textquotedblright{} is inherited: EtymDB takes it back through Middle Persian to an Indo-European word for \textquotedblleft{}ear\textquotedblright{}. \emph{bi} \textquotedblleft{}without\textquotedblright{} and \emph{shodan} \textquotedblleft{}to become\textquotedblright{} are native too.
\end{cousinweb}

\begin{grammarlens}[title={What you've built}]
One more word for telling, explaining and describing at B1.
\end{grammarlens}

\subsection*{Your turn}
\begin{itemize}
\raggedright
  \item \emph{Say it:} \emph{bihush shodan}
  \item \emph{Say it:} \emph{bihush shodan}, once more
  \item \emph{Say it:} a sentence of your own with \emph{bihush shodan}
  \item \emph{Recall:} say the Persian for a wound, then the Persian for bleeding, then say \emph{bihush shodan} again
\end{itemize}

\begin{tcolorbox}[breakable,colback=teal!4,colframe=teal!35!black,title={Before you move on}]
What does \fa{بیهوش} \fa{شدن} mean? (\textquotedblleft{}To faint, to pass out\textquotedblright{}.) Say it once more.
\end{tcolorbox}
\section[\texorpdfstring{\fa{به} \fa{هوش} \fa{آمدن} (be hush âmadan)}{به هوش آمدن (be hush âmadan)}]{\fa{به} \fa{هوش} \fa{آمدن} (be hush âmadan) — to come round, to regain consciousness}

\label{lesson:FA-C326-be-hush-amadan}



\begin{quote}
Before the new one: say the Persian for bleeding, then the Persian for to faint.
\end{quote}

\subsection*{You'll want to know: \fa{به} \fa{هوش} \fa{آمدن}}
\textbf{\fa{به} \fa{هوش} \fa{آمدن}} (\emph{be hush âmadan}) — \textquotedblleft{}to come round, to regain consciousness\textquotedblright{}.

\begin{quote}
\fa{او} \fa{بعد} \fa{از} \fa{ده} \fa{دقیقه} \fa{به} \fa{هوش} \fa{آمد}.
\end{quote}

\begin{quote}
\emph{U ba'd az dah daqiqe be hush âmad.}
\end{quote}

\begin{quote}
She came round after ten minutes.
\end{quote}

Literally \textquotedblleft{}to come to consciousness\textquotedblright{}. Present stem \emph{â-}: \emph{be hush miâyad}.

\begin{cousinweb}
\emph{hush} \textquotedblleft{}consciousness\textquotedblright{} is inherited from Middle Persian (EtymDB), and \emph{âmadan} \textquotedblleft{}to come\textquotedblright{} is native Persian too; EtymDB takes it back to the Indo-European root \emph{g\textsuperscript{w}em-} \textquotedblleft{}to step\textquotedblright{}.
\end{cousinweb}

\begin{grammarlens}[title={What you've built}]
One more word for telling, explaining and describing at B1.
\end{grammarlens}

\subsection*{Your turn}
\begin{itemize}
\raggedright
  \item \emph{Say it:} \emph{be hush âmadan}
  \item \emph{Say it:} \emph{be hush âmadan}, once more
  \item \emph{Say it:} a sentence of your own with \emph{be hush âmadan}
  \item \emph{Recall:} say the Persian for bleeding, then the Persian for to faint, then say \emph{be hush âmadan} again
\end{itemize}

\begin{tcolorbox}[breakable,colback=teal!4,colframe=teal!35!black,title={Before you move on}]
What does \fa{به} \fa{هوش} \fa{آمدن} mean? (\textquotedblleft{}To come round, to regain consciousness\textquotedblright{}.) Say it once more.
\end{tcolorbox}
